\section{Related Work}
\label{sec:related}

We survey two lines of work on uncertainty quantification for LLM hallucination detection---entropy-based and consistency-based methods---and highlight the gap our work addresses.

\subsection{Entropy-Based Uncertainty Quantification}

Token-level entropy measures uncertainty in next-token prediction by computing the entropy of the softmax distribution over vocabulary. High entropy indicates the model is uncertain about which token to generate, potentially signaling lack of knowledge.

\citet{kadavath2022language} demonstrated that LLMs exhibit calibration properties, with confidence scores correlating with correctness on question-answering tasks. Their P(True) method directly queries models about answer correctness. \citet{kuhn2023semantic} introduced \emph{semantic entropy}, which clusters semantically equivalent answers before computing entropy, addressing the problem that multiple phrasings of the same answer inflate raw entropy estimates. They evaluated on natural language generation tasks including question answering, but did not compare against consistency-based methods.

\citet{malinin2018predictive} provided theoretical foundations for distinguishing epistemic and aleatoric uncertainty in neural networks using prior networks. More recently, \citet{xiong2023llmconfidence} surveyed confidence estimation methods for LLMs, categorizing approaches by whether they require access to logits or only outputs.

\textbf{Limitation:} Entropy methods have primarily been evaluated on NLG tasks or with proprietary models, without systematic comparison to consistency-based approaches on factuality benchmarks.

\subsection{Consistency-Based Detection}

An orthogonal approach measures hallucination through response consistency across multiple generations. If a model generates semantically different answers to the same question across independent samples, this suggests uncertainty or fabrication.

\citet{manakul2023selfcheckgpt} introduced SelfCheckGPT, which generates multiple responses and computes consistency via BERTScore or embedding similarity. They demonstrated effectiveness on WikiBio biography generation, where inconsistent facts across samples indicate hallucination. \citet{wang2023selfconsistency} explored self-consistency in chain-of-thought reasoning, showing that majority voting across sampled reasoning chains improves accuracy.

\citet{chen2024contrastive} extended consistency methods with contrastive decoding, identifying hallucinations by comparing outputs from amateur and expert models. \citet{lin2022truthfulqa} created TruthfulQA specifically to evaluate truthfulness, showing that models often generate false but plausible-sounding answers.

\textbf{Limitation:} Consistency methods have been evaluated on different benchmarks (WikiBio, reasoning tasks) than entropy methods, making it impossible to determine whether the signals are redundant or complementary.

\subsection{The Gap: Missing Head-to-Head Comparison}

The two paradigms---entropy and consistency---have evolved independently:

\begin{table}[h]
\centering
\begin{tabular}{@{}llll@{}}
\toprule
\textbf{Method} & \textbf{Benchmark} & \textbf{Model} & \textbf{Direct Comparison} \\
\midrule
Semantic Entropy \citep{kuhn2023semantic} & NLG tasks & Various & No consistency baseline \\
SelfCheckGPT \citep{manakul2023selfcheckgpt} & WikiBio & GPT-3 & No entropy baseline \\
P(True) \citep{kadavath2022language} & QA tasks & Proprietary & No consistency baseline \\
\bottomrule
\end{tabular}
\caption{Prior work has not compared both methods on the same benchmark with the same model.}
\label{tab:gap}
\end{table}

No prior work has: (1) compared both methods on the \textbf{same factuality benchmark} (e.g., TruthfulQA), (2) using the \textbf{same model} (e.g., LLaMA-2-7B), (3) with analysis of whether signals are \textbf{redundant or orthogonal}.

This gap prevents practitioners from making informed choices about which uncertainty signal to use---or whether combining them provides additional value.

\subsection{Our Positioning}

We address this gap by conducting the first systematic comparison of token entropy and N-sample consistency on TruthfulQA with LLaMA-2-7B. Beyond comparing raw detection performance, we analyze the \emph{correlation} between signals and the \emph{complementary predictive value} on cases where methods disagree. This determines whether a hybrid approach is theoretically justified before investing in its implementation.
